\textbf{1. Kiến thức cần nhớ}
\begin{itemize}
\item Với $\vec a,\vec b\ne\vec0$, góc $(\vec a,\vec b)=\alpha\in[0^\circ;180^\circ]$ được xác định khi đặt hai véctơ chung gốc. Định nghĩa: $\boxed{\vec a\cdot\vec b=|\vec a||\vec b|\cos\alpha}$. Quy ước: $\vec a\cdot\vec0=\vec0\cdot\vec a=0$.
\item Tính chất: $\vec a\cdot\vec b=\vec b\cdot\vec a$; $\vec a\cdot(\vec b+\vec c)=\vec a\cdot\vec b+\vec a\cdot\vec c$; $(k\vec a)\cdot\vec b=k(\vec a\cdot\vec b)$.
\item $\vec a^{\,2}=\vec a\cdot\vec a=|\vec a|^2$; $(\vec a\pm\vec b)^2=|\vec a|^2\pm2\vec a\cdot\vec b+|\vec b|^2$; $(\vec a+\vec b)\cdot(\vec a-\vec b)=|\vec a|^2-|\vec b|^2$.
\item Với $\vec a,\vec b\ne\vec0$: $\boxed{\cos(\vec a,\vec b)=\dfrac{\vec a\cdot\vec b}{|\vec a||\vec b|}}$; $\boxed{\vec a\perp\vec b\iff\vec a\cdot\vec b=0}$.
\end{itemize}

\textbf{2. Phương pháp giải các dạng}

%% DANG 0H5?4-1
\textbf{Dạng 1. Tích vô hướng, góc giữa 2 véctơ}
\begin{enumerate}
\item Đặt hai véctơ chung gốc; xác định độ dài và góc giữa chúng.
\item Tính $\vec a\cdot\vec b=|\vec a||\vec b|\cos\alpha$; nếu cần, phân tích véctơ rồi dùng tính phân phối.
\end{enumerate}
Lưu ý: Đảo chiều một véctơ làm góc thành $180^\circ-\alpha$ và tích vô hướng đổi dấu.

%% DANG 0H5?4-2
\textbf{Dạng 2. Tìm góc nhờ tích vô hướng}
\begin{enumerate}
\item Tính $\vec a\cdot\vec b$ và hai độ dài, có thể dùng các hằng đẳng thức.
\item Tính $\cos\alpha=\dfrac{\vec a\cdot\vec b}{|\vec a||\vec b|}$, suy ra $\alpha\in[0^\circ;180^\circ]$.
\end{enumerate}
Lưu ý: Kiểm tra hai véctơ khác $\vec0$; góc giữa hai đường thẳng không vượt $90^\circ$.

%% DANG 0H5?4-3
\textbf{Dạng 3. Đẳng thức về tích vô hướng hoặc độ dài}
\begin{enumerate}
\item Dùng quy tắc ba điểm, trung điểm để biểu diễn các véctơ qua những véctơ chung.
\item Khai triển tích vô hướng, thay $\vec u^{\,2}=|\vec u|^2$ rồi rút gọn hai vế.
\item Với đẳng thức độ dài, có thể chứng minh bình phương hai vế bằng nhau nếu hai vế không âm.
\end{enumerate}
Lưu ý: $\vec a\cdot\vec b$ là một số; không dùng phép khử như với tích các số.

%% DANG 0H5?4-4
\textbf{Dạng 4. Điều kiện vuông góc}
\begin{enumerate}
\item Chọn hai véctơ khác $\vec0$ có phương theo hai đường thẳng cần xét.
\item Biểu diễn, khai triển tích vô hướng; chứng minh hoặc giải điều kiện $\vec u\cdot\vec v=0$.
\end{enumerate}
Lưu ý: Loại giá trị làm véctơ bằng $\vec0$ hoặc hình suy biến.

%% DANG 0H5?4-7
\textbf{Dạng 5. Toán thực tế áp dụng tích vô hướng}
\begin{enumerate}
\item Mô hình hóa bằng véctơ; xác định độ dài, góc và thống nhất đơn vị.
\item Áp dụng tích vô hướng; chẳng hạn công của lực không đổi: $A=\vec F\cdot\vec s=Fs\cos\alpha$.
\item Tính toán, kết luận theo tình huống và ghi đơn vị.
\end{enumerate}
Lưu ý: $\alpha$ là góc giữa lực và độ dời; công có thể âm, dương hoặc bằng $0$.
